\documentclass[preview]{standalone}

\usepackage{amsmath}
\usepackage{amssymb}
\usepackage{stellar}
\usepackage{definitions}
\usepackage{bettelini}

\begin{document}

\id{linearalgebra-orthonormal-bases}
\genpage

\section{Orthonormal collections}

\begin{snippetdefinition}{orthonormal-collection-definition}{Orthonormal collection}
    Let \((V,\langle\cdot,\cdot\rangle)\) be an \innerproductspace. A finite
    collection
    \[
        \{e_1,\ldots,e_m\}\subseteq V
    \]
    is \emph{orthonormal} if
    \[
        \langle e_i,e_j\rangle=\delta_{ij}
    \]
    for every \(1\leq i,j\leq m\), where
    \[
        \delta_{ij}
        =
        \begin{cases}
            1 & \text{if } i=j,\\
            0 & \text{if } i\neq j.
        \end{cases}
    \]
    Equivalently, every \(e_i\) has norm \(1\), and
    \[
        e_i\perp e_j
    \]
    whenever \(i\neq j\).
\end{snippetdefinition}

\begin{snippetexample}{orthonormal-collections-examples}{Orthonormal collections}
    In \(\mathbb{F}^n\) with the canonical \innerproduct, the canonical
    \basis
    \[
        \{e_1,\ldots,e_n\}
    \]
    is \orthonormal.

    In \(\realnumbers^3\), the collection
    \[
        \left\{
            \left(\frac{1}{\sqrt{3}},\frac{1}{\sqrt{3}},\frac{1}{\sqrt{3}}\right),
            \left(\frac{1}{\sqrt{2}},-\frac{1}{\sqrt{2}},0\right)
        \right\}
    \]
    is \orthonormal. The collection
    \[
        \left\{
            \left(\frac{1}{\sqrt{3}},\frac{1}{\sqrt{3}},\frac{1}{\sqrt{3}}\right),
            \left(\frac{1}{\sqrt{2}},-\frac{1}{\sqrt{2}},0\right),
            \left(\frac{1}{\sqrt{6}},\frac{1}{\sqrt{6}},-\frac{2}{\sqrt{6}}\right)
        \right\}
    \]
    is also \orthonormal in \(\realnumbers^3\).
\end{snippetexample}

\begin{snippetproposition}{orthonormal-collection-norm-formula}{Norm of a linear combination of orthonormal vectors}
    Let \((V,\langle\cdot,\cdot\rangle)\) be an \innerproductspace over
    \(\mathbb{F}\in\{\realnumbers,\complexnumbers\}\), and let
    \(\{e_1,\ldots,e_m\}\) be an \orthonormal collection in \(V\). Then, for
    every \(a_1,\ldots,a_m\in\mathbb{F}\),
    \[
        \left\|
        a_1e_1+\cdots+a_me_m
        \right\|^2
        =
        |a_1|^2+\cdots+|a_m|^2.
    \]
    In particular, \(\{e_1,\ldots,e_m\}\) is \linearlyindependent.
\end{snippetproposition}

\begin{snippetproof}{orthonormal-collection-norm-formula-proof}{orthonormal-collection-norm-formula}{Norm of a linear combination of orthonormal vectors}
    Since \(\{e_1,\ldots,e_m\}\) is \orthonormal,
    \[
        \langle e_i,e_j\rangle=\delta_{ij}.
    \]
    Therefore
    \[
        \left\|
        \sum_{i=1}^{m}a_ie_i
        \right\|^2
        =
        \left\langle
        \sum_{i=1}^{m}a_ie_i,
        \sum_{j=1}^{m}a_je_j
        \right\rangle
        =
        \sum_{i=1}^{m}\sum_{j=1}^{m}
        a_i\overline{a_j}\langle e_i,e_j\rangle
        =
        \sum_{i=1}^{m}|a_i|^2.
    \]
    If
    \[
        a_1e_1+\cdots+a_me_m=0,
    \]
    then the formula gives
    \[
        |a_1|^2+\cdots+|a_m|^2=0.
    \]
    Each term is nonnegative, so \(a_1=\cdots=a_m=0\). Hence the collection
    is \linearlyindependent.
\end{snippetproof}

\section{Orthonormal bases}

\begin{snippetdefinition}{orthonormal-basis-definition}{Orthonormal basis}
    Let \((V,\langle\cdot,\cdot\rangle)\) be an \innerproductspace. An
    \emph{orthonormal basis} of \(V\) is an \orthonormal collection
    \[
        \{e_1,\ldots,e_n\}\subseteq V
    \]
    that is a \basis of \(V\).
\end{snippetdefinition}

\begin{snippetexample}{orthonormal-bases-examples}{Orthonormal bases}
    The canonical \basis is an \orthonormalbasis of \(\mathbb{F}^n\) with
    the canonical \innerproduct.

    In \(\realnumbers^2\), the collection
    \[
        \left\{
            \left(\frac{1}{\sqrt{2}},\frac{1}{\sqrt{2}}\right),
            \left(\frac{1}{\sqrt{2}},-\frac{1}{\sqrt{2}}\right)
        \right\}
    \]
    is an \orthonormalbasis.

    In \(\realnumbers^4\), the collection
    \[
        \left\{
            \left(\frac{1}{2},\frac{1}{2},\frac{1}{2},\frac{1}{2}\right),
            \left(\frac{1}{2},\frac{1}{2},-\frac{1}{2},-\frac{1}{2}\right),
            \left(\frac{1}{2},-\frac{1}{2},-\frac{1}{2},\frac{1}{2}\right),
            \left(-\frac{1}{2},\frac{1}{2},-\frac{1}{2},\frac{1}{2}\right)
        \right\}
    \]
    is an \orthonormalbasis with respect to the canonical \innerproduct.
\end{snippetexample}

\begin{snippetproposition}{orthonormal-basis-associated-matrix}{Associated matrix of an orthonormal basis}
    Let \((V,\langle\cdot,\cdot\rangle)\) be a finite-dimensional
    \innerproductspace, and let
    \[
        \mathcal{B}=\{e_1,\ldots,e_n\}
    \]
    be a \basis of \(V\). Then \(\mathcal{B}\) is an \orthonormalbasis if
    and only if the \innerproductmatrix of \(\langle\cdot,\cdot\rangle\) in
    \(\mathcal{B}\) is
    \[
        \identmatrix{n}.
    \]
\end{snippetproposition}

\begin{snippetproof}{orthonormal-basis-associated-matrix-proof}{orthonormal-basis-associated-matrix}{Associated matrix of an orthonormal basis}
    The \innerproductmatrix in the \basis \(\mathcal{B}\) is
    \(S_{\mathcal{B}}=(s_{ij})\), where
    \[
        s_{ij}=\langle e_i,e_j\rangle.
    \]
    Thus \(S_{\mathcal{B}}=\identmatrix{n}\) if and only if
    \[
        \langle e_i,e_j\rangle=\delta_{ij}
    \]
    for every \(i,j\), which is exactly the condition that
    \(\mathcal{B}\) is \orthonormal.
\end{snippetproof}

\section{Coordinates}

\begin{snippetproposition}{parseval-identity}{Parseval identity}
    Let \((V,\langle\cdot,\cdot\rangle)\) be an \innerproductspace over
    \(\mathbb{F}\in\{\realnumbers,\complexnumbers\}\), and let
    \[
        \mathcal{B}=\{e_1,\ldots,e_n\}
    \]
    be an \orthonormalbasis of \(V\). Then, for every \(v\in V\),
    \[
        v=
        \langle v,e_1\rangle e_1
        +\cdots+
        \langle v,e_n\rangle e_n.
    \]
    Moreover,
    \[
        \|v\|^2
        =
        |\langle v,e_1\rangle|^2
        +\cdots+
        |\langle v,e_n\rangle|^2.
    \]
\end{snippetproposition}

\begin{snippetproof}{parseval-identity-proof}{parseval-identity}{Parseval identity}
    Since \(\mathcal{B}\) is a \basis, there exist scalars
    \(a_1,\ldots,a_n\in\mathbb{F}\) such that
    \[
        v=a_1e_1+\cdots+a_ne_n.
    \]
    For every \(j\),
    \[
        \langle v,e_j\rangle
        =
        a_1\langle e_1,e_j\rangle+\cdots+a_n\langle e_n,e_j\rangle
        =
        a_j.
    \]
    Hence
    \[
        v=
        \langle v,e_1\rangle e_1
        +\cdots+
        \langle v,e_n\rangle e_n.
    \]
    Applying the norm formula for \orthonormal collections gives
    \[
        \|v\|^2
        =
        |\langle v,e_1\rangle|^2
        +\cdots+
        |\langle v,e_n\rangle|^2.
    \]
\end{snippetproof}

\begin{snippetdefinition}{orthonormal-basis-fourier-coefficients-definition}{Fourier coefficients with respect to an orthonormal basis}
    Let \((V,\langle\cdot,\cdot\rangle)\) be an \innerproductspace over
    \(\mathbb{F}\in\{\realnumbers,\complexnumbers\}\), let
    \[
        \mathcal{B}=\{e_1,\ldots,e_n\}
    \]
    be an \orthonormalbasis of \(V\), and let \(v\in V\). The scalars
    \[
        \langle v,e_1\rangle,\ldots,\langle v,e_n\rangle
    \]
    are the \emph{Fourier coefficients of \(v\) with respect to
    \(\mathcal{B}\)}.
\end{snippetdefinition}

\section{Congruence}

\includesnpt{congruent-matrices-definition}

\begin{snippetproposition}{congruent-matrices-equivalence-relation}{Congruence is an equivalence relation}
    Let \(\mathbb{F}\in\{\realnumbers,\complexnumbers\}\). Congruence is an
    \equivrelation on \(\matrices_{n\times n}(\mathbb{F})\).
\end{snippetproposition}

\begin{snippetproof}{congruent-matrices-equivalence-relation-proof}{congruent-matrices-equivalence-relation}{Congruence is an equivalence relation}
    Congruence is reflexive because
    \[
        A=\identmatrix{n}^\transpose A\identmatrix{n}.
    \]
    It is symmetric because, if
    \[
        A=P^\transpose B\overline{P}
    \]
    for an invertible \matrix \(P\), then
    \[
        B=(P^{-1})^\transpose A\overline{P^{-1}}.
    \]
    It is transitive because, if
    \[
        A=P^\transpose B\overline{P}
        \qquad
        \text{and}
        \qquad
        B=Q^\transpose C\overline{Q},
    \]
    then
    \[
        A
        =
        P^\transpose Q^\transpose C\overline{Q}\,\overline{P}
        =
        (QP)^\transpose C\overline{QP}.
    \]
    Hence \(A\) is congruent to \(C\).
\end{snippetproof}

\begin{snippetproposition}{inner-product-matrices-are-congruent}{Associated matrices of the same inner product are congruent}
    Let \((V,\langle\cdot,\cdot\rangle)\) be a finite-dimensional
    \innerproductspace over \(\mathbb{F}\in\{\realnumbers,\complexnumbers\}\).
    Let \(\mathcal{B}\) and \(\mathcal{C}\) be \basis[bases] of \(V\), and
    let \(S_{\mathcal{B}}\) and \(S_{\mathcal{C}}\) be the
    \innerproductmatrix[associated matrices] of the same \innerproduct in
    the two bases. Then \(S_{\mathcal{B}}\) and \(S_{\mathcal{C}}\) are
    \congruentmatrices.
\end{snippetproposition}

\begin{snippetproof}{inner-product-matrices-are-congruent-proof}{inner-product-matrices-are-congruent}{Associated matrices of the same inner product are congruent}
    Let
    \[
        P=M(\identityfunc_V,\mathcal{B},\mathcal{C})
    \]
    be the change-of-basis \matrix from \(\mathcal{B}\) to
    \(\mathcal{C}\). Then, for every \(v\in V\),
    \[
        [v]_{\mathcal{C}}=P[v]_{\mathcal{B}}.
    \]
    For every \(u,v\in V\), the coordinate formula for the
    \innerproductmatrix gives
    \[
        \langle u,v\rangle
        =
        [u]_{\mathcal{B}}^\transpose
        S_{\mathcal{B}}
        \overline{[v]_{\mathcal{B}}}
    \]
    and also
    \[
        \langle u,v\rangle
        =
        [u]_{\mathcal{C}}^\transpose
        S_{\mathcal{C}}
        \overline{[v]_{\mathcal{C}}}.
    \]
    Substituting the coordinate-change formula yields
    \[
        \langle u,v\rangle
        =
        [u]_{\mathcal{B}}^\transpose
        P^\transpose S_{\mathcal{C}}\overline{P}
        \overline{[v]_{\mathcal{B}}}.
    \]
    Since this holds for every pair of coordinate vectors,
    \[
        S_{\mathcal{B}}=P^\transpose S_{\mathcal{C}}\overline{P}.
    \]
    Since \(P\) is invertible, this is exactly congruence.
\end{snippetproof}

\section{Gram-Schmidt}

\begin{snippettheorem}{gram-schmidt-theorem}{Gram-Schmidt process}
    Let \((V,\langle\cdot,\cdot\rangle)\) be an \innerproductspace over
    \(\mathbb{F}\in\{\realnumbers,\complexnumbers\}\), and let
    \[
        \{v_1,\ldots,v_n\}
    \]
    be a \linearlyindependent collection in \(V\). Define
    \[
        e_1=\frac{v_1}{\|v_1\|}.
    \]
    For \(j=2,\ldots,n\), define
    \[
        w_j
        =
        v_j-\sum_{i=1}^{j-1}\langle v_j,e_i\rangle e_i,
    \]
    and then
    \[
        e_j=\frac{w_j}{\|w_j\|}.
    \]
    Then \(\{e_1,\ldots,e_n\}\) is an \orthonormal collection and, for
    every \(1\leq j\leq n\),
    \[
        \linearspan(e_1,\ldots,e_j)
        =
        \linearspan(v_1,\ldots,v_j).
    \]
\end{snippettheorem}

\begin{snippetproof}{gram-schmidt-theorem-proof}{gram-schmidt-theorem}{Gram-Schmidt process}
    We argue by induction on \(j\). Since the starting collection is
    \linearlyindependent, \(v_1\neq0\). Hence \(e_1\) is well-defined and
    \[
        \|e_1\|=1.
    \]
    Also,
    \[
        \linearspan(e_1)=\linearspan(v_1).
    \]

    Suppose that \(e_1,\ldots,e_{j-1}\) have been constructed, are
    \orthonormal, and satisfy
    \[
        \linearspan(e_1,\ldots,e_{j-1})
        =
        \linearspan(v_1,\ldots,v_{j-1}).
    \]
    Since \(\{v_1,\ldots,v_n\}\) is \linearlyindependent,
    \[
        v_j\notin\linearspan(v_1,\ldots,v_{j-1})
        =
        \linearspan(e_1,\ldots,e_{j-1}).
    \]
    Therefore
    \[
        w_j=
        v_j-\sum_{i=1}^{j-1}\langle v_j,e_i\rangle e_i
    \]
    is nonzero, so \(e_j=w_j/\|w_j\|\) is well-defined.

    If \(k<j\), then
    \[
        \langle w_j,e_k\rangle
        =
        \left\langle
        v_j-\sum_{i=1}^{j-1}\langle v_j,e_i\rangle e_i,
        e_k
        \right\rangle
        =
        \langle v_j,e_k\rangle
        -
        \sum_{i=1}^{j-1}\langle v_j,e_i\rangle\langle e_i,e_k\rangle
        =
        0.
    \]
    Thus \(e_j\perp e_k\) for every \(k<j\), and by construction
    \[
        \|e_j\|=1.
    \]
    The enlarged collection is therefore \orthonormal.

    Since \(w_j\in\linearspan(v_1,\ldots,v_j)\), also
    \(e_j\in\linearspan(v_1,\ldots,v_j)\), and hence
    \[
        \linearspan(e_1,\ldots,e_j)
        \subseteq
        \linearspan(v_1,\ldots,v_j).
    \]
    Conversely,
    \[
        v_j=w_j+\sum_{i=1}^{j-1}\langle v_j,e_i\rangle e_i
        \in
        \linearspan(e_1,\ldots,e_j).
    \]
    Together with the induction hypothesis, this gives
    \[
        \linearspan(e_1,\ldots,e_j)
        =
        \linearspan(v_1,\ldots,v_j).
    \]
    The induction proves the result for every \(j\).
\end{snippetproof}

\begin{snippetcorollary}{orthonormal-basis-existence-corollary}{Existence of orthonormal bases}
    Every finite-dimensional \innerproductspace admits an
    \orthonormalbasis.
\end{snippetcorollary}

\begin{snippetproof}{orthonormal-basis-existence-corollary-proof}{orthonormal-basis-existence-corollary}{Existence of orthonormal bases}
    Let \(V\) be a finite-dimensional \innerproductspace. Choose any
    \basis of \(V\). Applying \gramschmidt to that \basis gives an
    \orthonormal collection with the same \linearspantext, hence an
    \orthonormalbasis of \(V\).
\end{snippetproof}

\begin{snippetcorollary}{orthonormal-completion-corollary}{Orthonormal completion}
    Let \(V\) be a finite-dimensional \innerproductspace. Every
    \orthonormal collection in \(V\) can be extended to an
    \orthonormalbasis of \(V\).
\end{snippetcorollary}

\begin{snippetproof}{orthonormal-completion-corollary-proof}{orthonormal-completion-corollary}{Orthonormal completion}
    Let \(\{e_1,\ldots,e_m\}\) be an \orthonormal collection in \(V\). It is
    \linearlyindependent, so by the
    \snippetref[linearly-independent-set-basis-expansion][basis extension theorem]
    it can be extended to a \basis
    \[
        \{e_1,\ldots,e_m,v_{m+1},\ldots,v_n\}
    \]
    of \(V\). Apply \gramschmidt to this basis. The first \(m\)
    vectors are unchanged because they are already \orthonormal. The
    remaining vectors are orthonormalized, and the resulting
    \orthonormalbasis contains \(e_1,\ldots,e_m\).
\end{snippetproof}

\begin{snippetexample}{gram-schmidt-polynomial-example}{Gram-Schmidt in \(\mathcal{P}_2(\mathbb{R})\)}
    Consider \(\mathcal{P}_2(\realnumbers)\) with the \innerproduct
    \[
        \langle p,q\rangle
        =
        p(0)q(0)+p(1)q(1)+p(-1)q(-1).
    \]
    Starting from the \basis
    \[
        \{1,x,x^2\},
    \]
    the first vector is normalized as
    \[
        \|1\|^2=\langle 1,1\rangle=3,
        \qquad
        e_1=\frac{1}{\sqrt{3}}.
    \]
    Since
    \[
        \langle x,e_1\rangle
        =
        \frac{1}{\sqrt{3}}(0+1-1)
        =
        0
    \]
    and
    \[
        \|x\|^2=0^2+1^2+(-1)^2=2,
    \]
    the second vector is
    \[
        e_2=\frac{x}{\sqrt{2}}.
    \]
    For the third vector,
    \[
        w_3=x^2-\langle x^2,e_1\rangle e_1-\langle x^2,e_2\rangle e_2.
    \]
    Now
    \[
        \langle x^2,e_1\rangle
        =
        \frac{1}{\sqrt{3}}(0^2+1^2+(-1)^2)
        =
        \frac{2}{\sqrt{3}},
    \]
    so
    \[
        \langle x^2,e_1\rangle e_1=\frac{2}{3}.
    \]
    Also,
    \[
        \langle x^2,e_2\rangle
        =
        \frac{1}{\sqrt{2}}\bigl(0^2\cdot0+1^2\cdot1+(-1)^2\cdot(-1)\bigr)
        =
        0.
    \]
    Hence
    \[
        w_3=x^2-\frac{2}{3}.
    \]
    Its norm satisfies
    \[
        \left\|x^2-\frac{2}{3}\right\|^2
        =
        \frac{4}{9}+\frac{1}{9}+\frac{1}{9}
        =
        \frac{2}{3},
    \]
    and therefore
    \[
        e_3
        =
        \sqrt{\frac{3}{2}}\left(x^2-\frac{2}{3}\right).
    \]
    Thus an \orthonormalbasis of \(\mathcal{P}_2(\realnumbers)\) is
    \[
        \left\{
            \frac{1}{\sqrt{3}},
            \frac{x}{\sqrt{2}},
            \sqrt{\frac{3}{2}}\left(x^2-\frac{2}{3}\right)
        \right\}.
    \]
\end{snippetexample}

\end{document}
